\chapter{\ploren{Podsumowanie i wnioski}{Conclusions}}

% Zastąp poniższe wskazówki treścią właściwą dla pracy. W podsumowaniu nie należy
% wprowadzać nowych metod, wyników ani źródeł, które nie były wcześniej omówione.

\begin{quote}
\itshape
\ploren
  {Przedstawiony układ jest propozycją. W krótszej pracy wszystkie elementy można połączyć w jeden spójny rozdział bez podrozdziałów. W pracy rozbudowanej podział pomaga jednoznacznie oddzielić ocenę celu, wkład własny, ograniczenia i kierunki dalszych prac. Tekst tej wskazówki należy usunąć z ostatecznej wersji pracy.}
  {The proposed structure is a recommendation. In a shorter thesis, all elements may be combined into one coherent chapter without subsections. In an extensive thesis, the division helps distinguish the evaluation of the objective, the author's contribution, limitations and directions for future work. This guidance should be removed from the final version of the thesis.}
\end{quote}

Podsumowanie powinno udzielać bezpośredniej odpowiedzi na problem postawiony w pracy. Nie jest powtórzeniem streszczenia ani spisu treści. Należy syntetycznie połączyć cel, zastosowaną metodę, najważniejsze rezultaty i wynikające z nich wnioski. Czytelnik powinien móc jednoznacznie ocenić, co zostało osiągnięte oraz jakie znaczenie ma wykonana praca.

\section{\ploren{Ocena realizacji celu pracy}{Evaluation of the thesis objective}}

Należy przypomnieć cel główny i jasno stwierdzić, czy został osiągnięty. Ocena powinna odwoływać się do konkretnych rezultatów oraz kryteriów określonych w rozdziale dotyczącym celu i zakresu. Nie wystarczy napisać, że „cel został zrealizowany”; trzeba krótko wskazać dowody potwierdzające tę ocenę.

Cele szczegółowe można omówić w tej samej kolejności, w jakiej zostały wcześniej sformułowane. Jeżeli któryś z nich zrealizowano tylko częściowo albo nie został osiągnięty, należy podać przyczynę i znaczenie tego ograniczenia dla celu głównego. Rzetelna ocena jest ważniejsza niż deklarowanie pełnego sukcesu niezgodnego z wynikami.

\section{\ploren{Najważniejsze wyniki i wnioski}{Principal results and conclusions}}

Należy wybrać kilka rezultatów, które najlepiej odpowiadają na pytania pracy. Wnioski powinny być formułowane na poziomie uzasadnionym przez dane i przeprowadzone badania. Jeśli rezultat dotyczy wyłącznie określonego stanowiska, zbioru danych, urządzenia lub zakresu parametrów, nie należy rozszerzać go bez podstaw na wszystkie możliwe warunki.

W pracy badawczej trzeba udzielić odpowiedzi na pytania badawcze albo podać wynik weryfikacji hipotez. W pracy projektowej należy ocenić spełnienie najważniejszych wymagań oraz użyteczność wykonanego rozwiązania. W pracy łączącej oba podejścia trzeba uwzględnić zarówno rezultat techniczny, jak i wnioski wynikające z jego badań.

\section{\ploren{Wkład własny autora}{Author's contribution}}

Należy syntetycznie wymienić najważniejsze elementy wykonane samodzielnie przez autora. Mogą to być opracowana metoda, projekt urządzenia, implementacja oprogramowania, model, stanowisko pomiarowe, procedura badawcza, zbiór danych, analiza wyników lub dokumentacja umożliwiająca odtworzenie pracy.

Wkład własny nie powinien być opisany jako lista zwykłych czynności, takich jak przeczytanie literatury lub uruchomienie gotowego programu. Trzeba wskazać rezultaty wymagające samodzielnych decyzji, pracy projektowej, implementacyjnej lub analitycznej oraz odróżnić je od wykorzystanych gotowych komponentów.

\section{\ploren{Znaczenie i możliwość wykorzystania rezultatów}{Significance and application of the results}}

Jeżeli wyniki mają zastosowanie praktyczne, należy wskazać możliwy sposób ich wykorzystania, warunki wdrożenia i potencjalnych odbiorców. W przypadku rezultatu poznawczego trzeba określić, jakie zrozumienie problemu zostało uzyskane i w czym uzupełnia ono stan wiedzy przedstawiony wcześniej.

Należy unikać deklaracji o znaczeniu rezultatów, których nie potwierdzono w pracy. Prototyp zweryfikowany laboratoryjnie nie jest jeszcze gotowym produktem, a skuteczność wykazana na ograniczonym zbiorze danych nie oznacza automatycznie poprawnego działania w środowisku rzeczywistym.

\section{\ploren{Ograniczenia pracy}{Limitations}}

Trzeba zebrać najważniejsze ograniczenia mające wpływ na interpretację i możliwość uogólnienia wyników. Nie ma potrzeby powtarzać całej dyskusji; należy wskazać te czynniki, które są istotne dla końcowej oceny rozwiązania. Mogą one dotyczyć zakresu badań, liczby prób, jakości danych, aparatury, uproszczeń modelu, platformy sprzętowej, czasu realizacji lub braku porównania z określonym wariantem.

Ograniczenia nie unieważniają pracy, jeżeli zostały rozpoznane, a wnioski sformułowano z odpowiednią ostrożnością. Powinny natomiast wyznaczać granice twierdzeń autora oraz prowadzić do konkretnych propozycji dalszej weryfikacji.

\section{\ploren{Kierunki dalszych prac}{Future work}}

Propozycje dalszych prac powinny wynikać z uzyskanych rezultatów i rozpoznanych ograniczeń. Zamiast ogólnego stwierdzenia, że rozwiązanie można rozwijać, należy wskazać konkretne działania: rozszerzenie zakresu danych, wykonanie badań w innych warunkach, zmianę konstrukcji, optymalizację algorytmu, porównanie z dodatkową metodą, przeprowadzenie testów długoterminowych albo przygotowanie wdrożenia.

Warto rozróżnić działania konieczne do pełnego potwierdzenia wniosków od opcjonalnych kierunków rozwoju funkcjonalnego. Nie należy przedstawiać jako przyszłej pracy elementu, którego wykonanie było niezbędne do osiągnięcia zadeklarowanego celu bieżącej pracy.

\section{\ploren{Wniosek końcowy}{Final conclusion}}

Rozdział należy zakończyć krótkim, jednoznacznym wnioskiem dotyczącym rezultatu całej pracy. Powinien on łączyć ocenę osiągnięcia celu z najważniejszym znaczeniem uzyskanych wyników, bez wprowadzania nowych informacji i bez przesadnego powtarzania wcześniejszych akapitów.
